% !TeX root = ../main.tex

\chapter{Methods}\label{chap:methods}

\section{TabPFN as the Surrogate}

\section{Expected Improvement from the Predictive Distribution}

\section{Acquisition Optimization and Trust Region}

\begin{figure}[t]
  \centering
  \resizebox{\textwidth}{!}{%
  \begin{tikzpicture}[
      item/.style={draw=#1, fill=#1!15, rounded corners=2pt, minimum height=7mm,
                   minimum width=34mm, align=center, font=\small},
      slot/.style={draw=#1, rounded corners=4pt, inner sep=3mm, line width=0.8pt},
      head/.style={font=\small\bfseries},
      arr/.style={-{Stealth[length=2.5mm]}, thick},
      lab/.style={font=\scriptsize, above, inner sep=2pt}
    ]
    % Surrogate column
    \node[item=surcol] (s1) at (0,-1)   {TabPFN};
    \node[item=surcol] (s2) at (0,-1.9) {GP (overlap)};
    \node[item=surcol] (s3) at (0,-2.8) {GP (trans.)};
    \node[item=surcol] (s4) at (0,-3.7) {GP (HED)};
    \node[slot=surcol, fit=(s1)(s2)(s3)(s4)] (S) {};
    \node[head, above=1mm of S] {Surrogate model};

    % Acquisition column (invisible pads give the frame the same height as the others)
    \node[minimum width=34mm, minimum height=7mm] (ap1) at (5.5,-1)   {};
    \node[minimum width=34mm, minimum height=7mm] (ap2) at (5.5,-3.7) {};
    \node[item=acqcol] (a1) at (5.5,-1.45) {EI};
    \node[item=acqcol] (a2) at (5.5,-2.35) {PI};
    \node[item=acqcol] (a3) at (5.5,-3.25) {LCB};
    \node[slot=acqcol, fit=(ap1)(ap2)] (A) {};
    \node[head, above=1mm of A] {Acquisition function};

    % Optimizer column
    \node[item=optcol] (o1) at (5.5*2,-1)   {Genetic algorithm};
    \node[item=optcol] (o2) at (5.5*2,-1.9) {Simulated annealing};
    \node[item=optcol] (o3) at (5.5*2,-2.8) {Hill climbing};
    \node[item=optcol] (o4) at (5.5*2,-3.7) {Bandit};
    \node[slot=optcol, fit=(o1)(o2)(o3)(o4)] (O) {};
    \node[head, above=1mm of O] {Acquisition optimizer};

    % Interfaces between components
    \draw[arr] (S.east) -- (A.west)
        node[lab, midway] {};
    \draw[arr] (A.east) -- (O.west)
        node[lab, midway] {};

    % Trust region
    \node[item=trcol, minimum width=40mm, below=6mm of O] (tr) {Trust region};
    \draw[arr] (tr.north) -- (O.south);
  \end{tikzpicture}}
  \caption{Modular components of BO. Each frame is a component whose options can be exchanged. The interfaces between components are the predictive distribution of the surrogate and the acquisition function, so in principle any option in one frame can be combined with the options in the others. The trust region restricts the search of the acquisition optimizer.}
  \label{fig:modular}
\end{figure}

\begin{figure}[t]
  \centering
  \resizebox{\textwidth}{!}{%
  \begin{tikzpicture}[
      font=\small,
      node distance=13mm,
      mod/.style={draw=#1, fill=#1!15, rounded corners=3pt, thick,
                  minimum width=33mm, minimum height=15mm, align=center},
      data/.style={draw=black!60, fill=black!5, rounded corners=3pt, thick,
                   minimum width=33mm, minimum height=15mm, align=center},
      arr/.style={-{Stealth[length=2.5mm]}, thick},
      lab/.style={font=\scriptsize, inner sep=2pt, align=center}
    ]
    % Row of stages
    \node[data] (in) {\textbf{Inputs}\\ $\mathcal{D}_n$, candidate $\boldsymbol{x}$};
    \node[mod=surcol, right=of in] (pre)
        {\textbf{Preprocessing}\\ scale continuous,\\ declare categorical};
    \node[mod=surcol, right=of pre, minimum width=42mm] (net)
        {\textbf{TabPFN forward pass}\\ context: $\mathcal{D}_n$\\ query: $\boldsymbol{x}$};
    \node[mod=surcol, right=of net] (ens)
        {\textbf{Ensemble}\\ average over\\ estimators};
    \node[mod=surcol, right=of ens] (bins)
        {\textbf{Predictive distribution}\\ bin probabilities\\ $p_k(\boldsymbol{x})$};
    \node[mod=acqcol, right=of bins] (ei)
        {\textbf{Expected improvement}\\ $\alpha_{\mathrm{EI}}(\boldsymbol{x})$ from $p_k$, $y_n^*$};

    \foreach \a/\b in {in/pre, pre/net, net/ens, ens/bins, bins/ei}
      \draw[arr] (\a) -- (\b);

    % To optimizer and back
    \node[mod=optcol, below=14mm of ei] (opt)
        {\textbf{Acquisition optimizer}\\ proposes new $\boldsymbol{x}$};
    \draw[arr] (ei) -- (opt) node[lab, midway, right] {$\alpha(\boldsymbol{x})$};
    \draw[arr] (opt.west) -| ($(in.south)+(0,-9mm)$) -- (in.south)
        node[lab, pos=0.15, below] {candidate $\boldsymbol{x}$ \\ (one forward pass per evaluation)};

    % Pretrained weights
    \node[draw=black!40, dashed, rounded corners=3pt, minimum width=42mm,
          minimum height=9mm, font=\scriptsize, text=black!60, above=9mm of net]
          (w) {pretrained weights (fixed)};
    \draw[arr, dashed, draw=black!40] (w) -- (net);
  \end{tikzpicture}}
  \caption{Computation of the acquisition value with TabPFN within one BO iteration. The observations serve as context and a candidate point as the query of a single forward pass, and the predictions of several estimators are averaged into a distribution over fixed bins, from which expected improvement is computed. The acquisition optimizer evaluates this pipeline once for each candidate point, and no parameters of TabPFN are updated.}
  \label{fig:tabpfn-pipeline}
\end{figure}

\begin{figure}[t]
  \centering
  \begin{tikzpicture}[font=\small]
    % bins
    \foreach \h [count=\i from 0] in {0.35,0.8,1.5,2.2,1.9,1.1,0.5,0.2} {
      \ifnum\i<4
        \fill[acqcol!55, draw=black!70] (\i*0.8,0) rectangle (\i*0.8+0.8,\h);
      \else
        \fill[surcol!25, draw=black!70] (\i*0.8,0) rectangle (\i*0.8+0.8,\h);
      \fi
    }
    \draw[thick, -{Stealth}] (-0.2,0) -- (6.9,0) node[right] {$y$};
    \draw[thick, dashed] (3.2,-0.15) -- (3.2,2.7) node[above] {$y_n^*$};
    \node[font=\scriptsize] at (1.6,-0.55) {bins with improvement};
    \node[font=\scriptsize] at (4.8,-0.55) {bins without improvement};
    \node[font=\scriptsize, above] at (0.4,0.35) {$p_1$};
    \node[font=\scriptsize, above] at (1.2,0.8) {$p_2$};
    \node[font=\scriptsize, above] at (2.0,1.5) {$p_3$};
    \node[font=\scriptsize, above] at (2.8,2.2) {$p_4$};
  \end{tikzpicture}
  \caption{Predictive distribution of TabPFN over fixed bins at a candidate point. The expected improvement is the sum, over the bins below the best observation $y_n^*$ (orange), of the bin probability $p_k$ multiplied by the expected improvement within the bin.}
  \label{fig:ei-bins}
\end{figure}

\begin{figure}[t]
  \centering
  \resizebox{\textwidth}{!}{%
  \begin{tikzpicture}[
      font=\small,
      node distance=12mm and 16mm,
      mod/.style={draw=#1, fill=#1!15, rounded corners=3pt, thick,
                  minimum width=36mm, minimum height=13mm, align=center},
      data/.style={draw=black!60, fill=black!5, rounded corners=3pt, thick,
                   minimum width=36mm, minimum height=13mm, align=center},
      arr/.style={-{Stealth[length=2.5mm]}, thick},
      lab/.style={font=\scriptsize, inner sep=2pt}
    ]
    % Top row
    \node[data] (init) {\textbf{Initial design}\\ $\boldsymbol{x}_1,\dots,\boldsymbol{x}_{n_0}$};
    \node[data, right=of init] (data) {\textbf{Observations}\\ $\mathcal{D}_n$};
    \node[mod=surcol, right=of data] (sur)
        {\textbf{Surrogate model}\\ $p(y \mid \boldsymbol{x}, \mathcal{D}_n)$};
    \node[mod=acqcol, right=of sur] (acq)
        {\textbf{Acquisition function}\\ $\alpha(\boldsymbol{x}; \mathcal{D}_n)$};

    % Bottom row
    \node[mod=optcol, below=of acq] (opt)
        {\textbf{Acquisition optimizer}\\ $\boldsymbol{x}_{n+1}=\arg\max_{\boldsymbol{x}}\alpha$};
    \node[mod=trcol, left=of opt] (tr) {\textbf{Trust region}\\ restricts the search};
    \node[data, below=of data] (obj)
        {\textbf{Objective}\\ $y_{n+1}=f(\boldsymbol{x}_{n+1})+\epsilon$};

    % Arrows
    \draw[arr] (init) -- (data);
    \draw[arr] (data) -- (sur);
    \draw[arr] (sur) -- (acq);
    \draw[arr] (acq) -- (opt);
    \draw[arr] (tr) -- (opt);
    \draw[arr] (opt.south) -- ++(0,-8mm) -| (obj.south)
        node[lab, pos=0.25, below] {$\boldsymbol{x}_{n+1}$};
    \draw[arr] (obj.north) -- (data.south)
        node[lab, midway, right] {$(\boldsymbol{x}_{n+1}, y_{n+1})$};
  \end{tikzpicture}}
  \caption{The BO loop. The surrogate, the acquisition function, and the acquisition optimizer are modular components, and the trust region restricts the search of the optimizer. The initial design is evaluated on $f$ to form the first observations.}
  \label{fig:bo-loop}
\end{figure}

\begin{figure}[t]
  \centering
  \resizebox{\textwidth}{!}{%
  \begin{tikzpicture}[
      font=\small,
      node distance=22mm,
      blk/.style={draw=#1, fill=#1!15, rounded corners=3pt, thick,
                  minimum width=36mm, minimum height=14mm, align=center},
      dat/.style={draw=black!60, fill=black!5, rounded corners=3pt, thick,
                  minimum width=36mm, minimum height=11mm, align=center},
      arr/.style={-{Stealth[length=2.5mm]}, thick},
      lab/.style={font=\scriptsize, align=center, inner sep=2pt}
    ]
    % Inputs
    \node[dat] (D) at (0,0.9)  {Observations\\ $\mathcal{D}_n$};
    \node[dat] (X) at (0,-0.9) {Candidate\\ $\boldsymbol{x}$};

    % Predictive distribution
    \node[blk=surcol] (P) at (7.5,0)
        {Predictive distribution\\ $p(y \mid \boldsymbol{x}, \mathcal{D}_n)$};

    % Acquisition value
    \node[blk=acqcol] (A) at (15,0)
        {Acquisition value\\ $\alpha(\boldsymbol{x}; \mathcal{D}_n)$};

    % Optimizer
    \node[blk=optcol] (O) at (22.5,0) {Acquisition\\ optimizer};

    % Merge inputs
    \coordinate (m) at (2.7,0);
    \draw[thick] (D.east) -| (m);
    \draw[thick] (X.east) -| (m);
    \draw[arr] (m) -- (P.west)
        node[lab, midway, above] {forward pass\\ (TabPFN)};

    % Pretrained weights
    \node[draw=black!40, dashed, rounded corners=3pt, font=\scriptsize,
          text=black!60, minimum height=7mm] (W) at (5.1,2.2) {pretrained weights (fixed)};
    \draw[arr, dashed, black!40] (W.south) -- (5.1,0.35);

    % P -> A -> O
    \draw[arr] (P) -- (A)
        node[lab, midway, above] {expected\\ improvement\\ with $y_n^*$};
    \draw[arr] (A) -- (O);

    % Feedback of new candidates
    \draw[arr] (O.south) -- ++(0,-12mm) -| (X.south)
        node[lab, pos=0.25, below] {new candidate $\boldsymbol{x}$ \\ (one forward pass per candidate)};
  \end{tikzpicture}}
  \caption{Computation of the acquisition value with TabPFN within one BO iteration. The observations serve as context and a candidate point as the query of a single forward pass of the pretrained network, which returns a predictive distribution over fixed bins. Expected improvement is computed from this distribution and the best observation $y_n^*$, and the acquisition optimizer proposes new candidate points until the search ends.}
  \label{fig:tabpfn-pipeline}
\end{figure}

\begin{figure}[t]
  \centering
  \resizebox{\textwidth}{!}{%
  \begin{tikzpicture}[
      font=\small,
      node distance=14mm and 14mm,
      mod/.style={draw=#1, fill=#1!15, rounded corners=3pt, thick,
                  minimum width=30mm, minimum height=12mm, align=center},
      data/.style={draw=black!60, fill=black!5, rounded corners=3pt, thick,
                   minimum width=30mm, minimum height=12mm, align=center},
      role/.style={font=\small\bfseries, inner sep=2pt},
      arr/.style={-{Stealth[length=2.5mm]}, thick},
      lab/.style={font=\scriptsize, inner sep=2pt}
    ]
    % Top row
    \node[data] (init) {Initial design\\ $\boldsymbol{x}_1,\dots,\boldsymbol{x}_{n_0}$};
    \node[data, right=of init] (data) {Observations\\ $\mathcal{D}_n$};
    \node[mod=surcol, right=of data] (sur) {\textbf{TabPFN}};
    \node[mod=acqcol, right=of sur] (acq) {\textbf{EI}};
    \node[role, above=1mm of sur] {Surrogate model};
    \node[role, above=1mm of acq] {Acquisition function};

    % Bottom row
    \node[data, below=of data] (obj)
        {Objective\\ $y_{n+1}=f(\boldsymbol{x}_{n+1})+\epsilon$};
    \node[mod=optcol, below=of acq] (opt) {\textbf{[GA / SA]}};
    \node[mod=trcol, right=of opt] (tr) {\textbf{Hamming + box}};
    \node[role, below=1mm of opt] {Acquisition optimizer};
    \node[role, below=1mm of tr] {Trust region};

    % Arrows
    \draw[arr] (init) -- (data);
    \draw[arr] (data) -- (sur);
    \draw[arr] (sur) -- (acq)
        node[lab, midway, above] {$p(y \mid \boldsymbol{x}, \mathcal{D}_n)$};
    \draw[arr] (acq) -- (opt)
        node[lab, midway, right] {$\alpha(\boldsymbol{x}; \mathcal{D}_n)$};
    \draw[arr] (tr) -- (opt);
    \draw[arr] (opt.west) -- (obj.east)
        node[lab, midway, above] {$\boldsymbol{x}_{n+1}$};
    \draw[arr] (obj.north) -- (data.south)
        node[lab, midway, right] {$(\boldsymbol{x}_{n+1}, y_{n+1})$};
  \end{tikzpicture}}
  \caption{The BO loop used in this thesis. The role of each component is shown outside its box and the choice made in this thesis inside. The trust region restricts the search of the acquisition optimizer. The initial design is evaluated on $f$ to form the first observations.}
  \label{fig:bo-loop-tabpfn}
\end{figure}